% A.84 reserved by the root.
\subsection{A full Tate equation coupling two bad places}

\begin{statement}{A.84}{Proved}
{The exact Tate logarithms recognize the parameter thirty-five}
Every noncuspidal Selmer section in A.14 with labeled integral
realization $E_{35}$ is the point section $s_{35}$.
No square-multiplier or full-$H$ hypothesis is imposed.
The proof retains one global multiplier and compares its
coefficient-five squareclass at the distinct local primes
five and seventeen.
\end{statement}
\begin{proof}
Fix the global integral intertwiners in the direction
$g_\ell:T_\ell(E_{35})\to T_\ell(E_U)|_s$ and their
Weil multipliers $\Delta_\ell=\det g_\ell$ as in A.82.
The potentially multiplicative primes are $5,7,17$.
Their labels are respectively zero, zero and one by the
labeled cyclotomic-line comparison.
At two the realized curve is potentially good, since
$v_2(35)=0$, $v_2(1-35)=1$ and $v_2(j)=6$.
There is no dyadic cusp.

We determine its valuation pair before using any
determinant restriction.  The squareclasses in A.16
give even valuation parity for $A_2$ and odd parity for
$B_2$.  Among the potentially good dyadic pairs this
allows only $(0,1)$ and $(0,3)$.
The second pair would make $x_2$ one modulo eight and
therefore a square, while its prescribed squareclass
is $35\equiv3\pmod8$.  Thus the pair is exactly
$(0,1)$.
Signs and prime valuations now give the complete global
Kummer expansions
\[
                    A=5^\beta7^\gamma,
                  \qquad B=-2\cdot17^\alpha,
       \qquad\alpha,\beta,\gamma\in\widehat\Z^\times.
\tag{A.84.1}
\]
The unit assertion follows from the full integral Tate
valuation comparison at the three odd bad primes,
each of whose realized contact orders is one.
It also holds at any such prime which is ordinary,
where A.78 gives its exact parameter and contact order.
The common-multiplier calculation (A.82.7) gives,
for every coefficient $\ell$,
\[
 \alpha_\ell=\frac{u_{17,\ell}^2}{\Delta_\ell},
 \qquad
 \beta_\ell=\frac{u_{5,\ell}^2}{\Delta_\ell},
 \qquad
 \gamma_\ell=\frac{u_{7,\ell}^2}{\Delta_\ell},
       \qquad u_{q,\ell}\in\Z_\ell^\times.
\tag{A.84.2}
\]

Suppose that $s$ is not a point section.
The one-label theorem A.37 then requires both cusp
labels zero and one.  The only possible label-one
prime is seventeen, so $s_{17}$ is a cusp at one.
We use its full coefficient-seventeen Tate extension.
The actual Tate parameter of the realized curve is
\[
              q_{17}=\frac{(-34)^2}{256}F(-34)
                       =\frac{17^2}{64}F(-34),
\tag{A.84.3}
\]
where $F(t)\in1+t\Z[1/2][[t]]$ and
$F(t)=1+t+O(t^2)$ as in A.74.
Indeed, $w=1-t$ gives $w=-34$, the Legendre $j$ invariant
is unchanged, and the inverse $j$ series uniquely
specifies this ordinary Tate parameter.
The cusp-one Tate class is exactly $B^2/256$ by A.63.
Its valuation is $2\alpha_{17}$, so the diagonal ratio
in the comparison is $\alpha_{17}$.
Applying the logarithm with $\log_{17}(17)=0$ to the
full extension equality gives
\[
 -6\log_{17}(2)
      =\alpha_{17}
          \bigl(-6\log_{17}(2)+\log_{17}F(-34)\bigr).
\tag{A.84.4}
\]
Both representations are split at seventeen: the
realized curve has the two nodal slopes $\pm1$ at
contact one, as does its formal cusp representation.
No quadratic character has been discarded.

We justify the needed nome congruences uniformly.
For every odd $p$ and $w\in p\Z_p$, integrality away
from two gives $F(w)-1-w\in p^2\Z_p$.
In the logarithm of $1+(F(w)-1)$, every term of
degree $n\ge2$ has valuation at least
$n-v_p(n)\ge2$.
Consequently the exact congruence is
\[
                         \log_p F(w)\equiv w\pmod{p^2}.
\tag{A.84.5}
\]
Thus no unbounded denominators in the logarithm or
the formal nome have been ignored.
The Fermat quotient satisfies
$(2^{16}-1)/17\equiv13\pmod{17}$; hence
$\log_{17}(2)/17\equiv13/16=4\pmod{17}$.
After dividing (A.84.4) by seventeen, (A.84.5) makes
its numerator $-6\cdot4=10$ and its denominator
$10-2=8$ modulo seventeen.  In particular the latter
has exact valuation one, so division is legitimate and
\[
                           \alpha_{17}\equiv10/8=14
                                            \pmod{17}.
\tag{A.84.6}
\]
Fourteen is a nonsquare modulo seventeen.
Equation (A.84.2) therefore makes $\Delta_{17}$ a
nonsquare as well.
An ordinary restriction at any potentially multiplicative
place would make the entire multiplier a square by
A.34.  It follows that five and seven, as well as
seventeen, must all be cuspidal under the supposition.

At five the label is zero, so the regular coordinate
$B_5$ equals one.  Its coefficient-five logarithm in
(A.84.1) is
\[
                 \log_5(2)+\alpha_5\log_5(17)=0.
\tag{A.84.7}
\]
The exact small congruences are
\[
       \log_5(2)\equiv10,\qquad
       \log_5(17)\equiv5,\qquad
       \log_5(7)\equiv0\pmod{25}.
\tag{A.84.8}
\]
They follow by applying
$\log_5(a)\equiv(a^4-1)/4\pmod{25}$ to
$2^4=16$, $17^4\equiv21\pmod{25}$ and
$7^4=2401\equiv1\pmod{25}$.
Thus both first logarithms have valuation one, and
(A.84.7) gives
\[
                                    \alpha_5\equiv3\pmod5.
\tag{A.84.9}
\]

The full Tate extension at this same local prime
provides an independent comparison of $\beta_5$.
Its realized parameter and logarithm are
\[
 q_5=\frac{35^2}{256}F(35),
 \qquad
 \log_5(q_5)=2\log_5(7)-8\log_5(2)+\log_5F(35),
\tag{A.84.10}
\]
where $\log_5(5)=0$.
The curve and cusp-zero splitting character is that
of $-1$, which is trivial over $\Q_5$.
The cusp class $A^2/256$ has valuation $2\beta_5$;
the realized parameter has valuation two.
The actual coefficient-five extension equality is therefore
\[
        2\gamma_5\log_5(7)-8\log_5(2)
                             =\beta_5\log_5(q_5).
\tag{A.84.11}
\]
By (A.84.5), $\log_5F(35)\equiv35\pmod{25}$.
Equations (A.84.8) and (A.84.10) give
\[
                         \log_5(q_5)\equiv5\pmod{25}.
\tag{A.84.12}
\]
Since $\gamma_5$ is integral, its term in (A.84.11)
vanishes modulo twenty-five.  Division of that equality
by five now gives
\[
                                      \beta_5\equiv4\pmod5.
\tag{A.84.13}
\]
But (A.84.2), at the distinct local primes seventeen
and five with the same coefficient five, requires
\[
                 \frac{\alpha_5}{\beta_5}
                       =\left(\frac{u_{17,5}}{u_{5,5}}\right)^2
                         \in(\Z_5^\times)^2.
\tag{A.84.14}
\]
Its residue from (A.84.9) and (A.84.13) is
$3/4=2\pmod5$, a nonsquare.
This contradiction proves that $s$ is a rational point
section.  Its exact parameter is $35$ by the labeled
primitive-isogeny argument of A.54.
\end{proof}

\paragraph{What the full extension equation contributes.}
At the label-one prime seventeen, the regular coordinate
alone gives the unit-square form with coefficients
$\log_{17}(5)$ and $\log_{17}(7)$.
Their normalized Fermat quotients are twelve and three,
so its required square ratio is $-12/3=13\pmod{17}$.
Thirteen is a square, and the simple-root Hensel argument
of A.82 gives an actual unit zero of this local quadratic
form.  Thus the regular-coordinate quadratic test has
an explicit surviving local solution here.
Equation (A.84.4) uses the additional full Tate class;
it fixes the multiplier squareclass at coefficient
seventeen and thereby makes the coefficient-five
comparison (A.84.14) available at both bad places.
Only the latter full calculation proves recognition.

The exact residue, logarithm and inverse-nome checks
are retained in
\path{code/legendre_parameter_35_certificate.py}
and \path{code/legendre_parameter_35_certificate_output.txt}.
The finite checks do not replace the common-multiplier
and full extension maps proved above.
